\documentclass{article}
\usepackage[a4paper, left=1in, top=1in, right=1in, bottom=1in]{geometry}
\usepackage{amsmath}
\usepackage{graphicx}
\usepackage[style=numeric]{biblatex} % Set citation/bibliography style
\addbibresource{references.bib}
\title{ICS 691G: An Exploration of Hidden Markov Models In Regards to Stability, Underflow, and External Memory}
\author{Hunter Von Tungeln,}
\date{October 6, 2026}
\begin{document}
\maketitle
\section{Background}
Hidden Markov Models (HMMs) are an important tool in biological sequence alignment \cite{rabiner1989tutorial}. Two such algorithms that utilize HMMs are the Forward and Viterbi algorithms.  These algorithms are very memory heavy, requiring storing many floating points and performing many floating-point operations. In addition, since these algorithms work with mathematical probabilities, they deal with very small numerical values, which can lead to underflow in computer systems. To mitigate underflow, these algorithms perform their calculations in ``log-space", which means they convert the probabilities they work with into their natural logarithms. All this being said, memory usage is a big concern for these algorithms, and ways to improve memory usage would make these algorithms more effective. 
\section{Goal \& Objective}
The goal of this project is to empirically analyze different implementations of the Forward and Viterbi HMM algorithms.  This project seeks to explore the different strategies used to combat numerical underflow across different floating-point precisions. This project also seeks to explore empirical behavior of the Forward and Viterbi algorithms within the external memory model, that is, measuring their total I/O operations and comparing the empirical behavior to theoretical results. 

\section{Scope \& Methodology}
This project will implement the Forward and Viterbi algorithms in C++. Different implementations will be studied to better understand combating underflow utilizing different floating point precisions (float, double, long double). The utilization of different floating point precisions is to compare at what point does underflow occur without any normalization methods, which will demonstrate the need for a form of scaling. Both algorithms will have a version implemented with only log-space calculations. Both algorithms will also have another version using a hybrid approach of Just-In-Time (JIT) scaling, where normal multiplication is used until a calculated value is within a threshold of underflow,  which will then rescale the calculated probabilities. These problems can be formulated as follows: \\ \\
\textbf{Hypothesis:} If the Forward and Viterbi Algorithms are applied to the Human Reference Genome, using three different floating point precisions (float, double, long double), then the lower precision data types will cause an underflow faster relative to the higher precision data types. \\\\
\textbf{Hypothesis:} If the Forward and Viterbi Algorithms are applied to the Human Reference Genome, using either log-space or Just-In-Time scaling, then the implementations utilizing Just-In-Time scaling will better utilize memory due to less total computation being done in log space. \\\\
\textbf{Hypothesis:} If the Forward and Viterbi Algorithms are applied to the Human Reference Genome, then they will exhibit I/O behavior synonymous of the their theoretical external memory model bounds. \\\\

The performance of these algorithms will be quantified in different ways. The CPU time of the log-space and JIT scaling implementations will be measured and compared (using something like perf). Lower time is better. The I/O and memory overhead incurred by both of these algorithms will be measured (using something like Valgrind). Lower I/O and memory use is better. The number of I/O operations will be plotted for different sequence lengths, to compare to the theoretical upper bounds in the external memory model. Currently, the plan is to benchmark these implementations on the Human Reference Genome, which can be found in the UCSC Genome Browser here \cite {Kent2002}\cite{Miga697-707}. 
\section{Timeline}
\begin{description}
    \item[\textbf{Oct 6:}] Submit formal project proposal.
    \item[\textbf{Oct 7 - Oct 31:}] Implement baseline Forward/Viterbi algorithms in C++ and their variants utilizing different floating point precisions
    \item[\textbf{Nov 1 - Nov 16:}] Implement Log-Space-only and JIT Rescaling variants.
    \item[\textbf{Nov 17 or Nov 19:}] Mid-term Checkpoint
    \item[\textbf{Nov 20 - Dec 8:}] Conduct empirical profiling (memory analysis, CPU benchmarking, precision underflow tests) on Koa.
    \item[\textbf{Dec 9 - Dec 18:}] Finalize the 10--15 page written report and prepare the final presentation. Submit code repository and report by the 5:00\,PM deadline on December 18.
\end{description}

\printbibliography

\end{document}